\chapter{Дослідження методу автоматичної ідентифікації музичних інструментів}
\label{chap:experiments}

У цьому розділі описано програмне забезпечення, розроблене для
дослідження методу автоматичної ідентифікації музичних інструментів,
математичне забезпечення якого викладено у розділі~\ref{chap:mathmodel},
а також наведено результати експериментального дослідження якості
запропонованого підходу.

% ------------------------------------------------------------------
\section{Засоби реалізації}
\label{sec:tools}

Програмне забезпечення реалізовано мовою Python~3. Для обробки
аудіосигналів та виділення мел-спектрограм використано бібліотеку
\texttt{librosa}~\cite{mcfee2015}, яка надає зручний інтерфейс для
обчислення STFT та мел-фільтрбанку. Побудову і навчання моделі
реалізовано з використанням фреймворку \texttt{PyTorch}. Чисельні
операції та роботу з масивами даних забезпечено бібліотекою
\texttt{NumPy}. Для обчислення метрик multi-label класифікації
використано бібліотеку \texttt{scikit-learn}.

Реалізація складається з таких основних модулів:
\begin{enumerate}[label=\asbuk*)]
    \item модуль попередньої обробки аудіосигналу та виділення ознак;
    \item модуль формування навчального набору даних;
    \item модуль визначення архітектури та навчання моделі;
    \item модуль оцінювання якості та візуалізації результатів.
\end{enumerate}

% ------------------------------------------------------------------
\section{Набір даних та підготовка}
\label{sec:dataset}

Для навчання та оцінювання методу використано датасет
IRMAS~\cite{irmas2012}, що містить близько 6000 тризекундних
аудіофрагментів з анотаціями для 11 класів інструментів: гітара,
електрогітара, орган, фортепіано, труба, скрипка, флейта, ударні,
контрабас, мандоліна та голос. Фрагменти отримано з реальних музичних
записів різних жанрів.

Оскільки один фрагмент може містити кілька інструментів одночасно,
мітки подаються у форматі бінарного вектора розмірності~11. Набір
даних розділено на навчальну та тестову вибірки у співвідношенні
$80/20$ із стратифікацією за розподілом класів.

Для розширення навчальної вибірки та підвищення стійкості моделі
застосовувалась аугментація даних: випадкові зсуви висоти тону в
межах $\pm 2$ напівтонів та зміна темпу в межах $\pm 10\%$.

% ------------------------------------------------------------------
\section{Алгоритм виділення ознак}
\label{sec:feature_extraction}

Для кожного аудіозапису виконувалась така послідовність операцій
обробки сигналу.

Аудіосигнал завантажувався з частотою дискретизації $f_s = 22050$~Гц
та нормалізувався за амплітудою. Мел-спектрограма обчислювалась з
параметрами: кількість мел-смуг $M = 128$, довжина вікна STFT
$L = 1024$, крок $H = 512$ відліків (близько 23~мс на фрейм).

Логарифмічне масштабування:
\begin{equation}
    \tilde{S}_t(m) = \log\!\left(S_t(m) + \epsilon\right),
    \quad \epsilon = 10^{-9},
    \label{eq:logmel}
\end{equation}
де $S_t(m)$ --- значення $m$-ї мел-смуги у момент часу $t$.

У результаті для кожного тризекундного фрагменту формувалась матриця
ознак $\mathbf{S} \in \mathbb{R}^{128 \times 130}$, що подавалась на
вхід моделі як одноканальне зображення.

% ------------------------------------------------------------------
\section{Архітектура та конфігурація моделі}
\label{sec:architecture}

Як модель multi-label класифікації використано гібридну архітектуру
CNN~+~BiLSTM. CNN-блок витягує тембральні ознаки зі спектрального
представлення, BiLSTM-блок моделює часовий контекст звучання.

Архітектура моделі включає:
\begin{enumerate}[label=\asbuk*)]
    \item блок CNN: два шари Conv2D з фільтрами $3 \times 3$, кількість
    карт 32 та 64 відповідно, активація ReLU; після кожного шару ---
    BatchNorm2D та MaxPool2D $(2 \times 2)$;
    \item шар Dropout з імовірністю $0.25$ після CNN-блоку;
    \item перетворення вихідного тензору CNN у послідовність фреймів;
    \item один шар BiLSTM з прихованою розмірністю $128$
    (по кожному напрямку);
    \item шар Dropout з імовірністю $0.3$ після BiLSTM;
    \item повнозв'язний вихідний шар розмірності $K = 11$
    (за кількістю класів інструментів);
    \item незалежна сигмоїдна активація для кожного класу.
\end{enumerate}

Навчання моделі здійснювалось з такими параметрами: функція втрат ---
бінарна крос-ентропія (BCELoss); оптимізатор --- Adam з початковою
швидкістю навчання $10^{-3}$; розмір батчу --- $64$ фрагменти;
кількість епох --- $50$. Для уникнення перенавчання застосовувався
ранній зупин з терпінням 10 епох за значенням F1-macro на
валідаційній вибірці.

% ------------------------------------------------------------------
\section{Схема тестування та метрики оцінювання}
\label{sec:testing}

% TODO

% ------------------------------------------------------------------
\section{Результати дослідження}
\label{sec:results}

% TODO

\subsection{Динаміка навчання моделі}
\label{subsec:training_dynamics}

% TODO

\subsection{Порівняння з базовими методами}
\label{subsec:comparison}

% TODO

\subsection{Аналіз результатів по класах}
\label{subsec:class_analysis}

% TODO

% ------------------------------------------------------------------
\section{Висновки до розділу}
\label{sec:ch4_conclusions}

% TODO
